\section{Energy--momentum tensor with dimensional regularization}
\label{sec:2}
The description of the energy--momentum tensor in gauge
theory~\cite{Freedman:1974gs,Joglekar:1975jm} is particularly simple with the
dimensional regularization.\footnote{Ref.~\cite{Collins:1984xc} is a very nice
exposition of the dimensional regularization.} This is because this
regularization manifestly preserves the (vectorial) gauge symmetry and the
translational invariance. Thus, in this section, we briefly recapitulate basic
facts concerning the energy--momentum tensor on the basis of the dimensional
regularization.

The action of the system under consideration in a $D$~dimensional Euclidean
space is given by
\begin{equation}
   S=\frac{1}{4g_0^2}\int\mathrm{d}^Dx\,F_{\mu\nu}^a(x)F_{\mu\nu}^a(x)
   +\int\mathrm{d}^Dx\,\Bar{\psi}(x)(\Slash{D}+m_0)\psi(x),
\label{eq:(2.1)}
\end{equation}
where $g_0$ and~$m_0$ are bare gauge coupling and the mass parameter,
respectively. The field strength is defined by
\begin{equation}
   F_{\mu\nu}(x)
   =\partial_\mu A_\nu(x)-\partial_\nu A_\mu(x)+[A_\mu(x),A_\nu(x)],
\label{eq:(2.2)}
\end{equation}
for $A_\mu(x)=A_\mu^a(x)T^a$ and~$F_{\mu\nu}(x)=F_{\mu\nu}^a(x)T^a$, and the
covariant derivative on the fermion is
\begin{equation}
   D_\mu=\partial_\mu+A_\mu.
\label{eq:(2.3)}
\end{equation}
Here, and in what follows, the summation over $N_{\mathrm{f}}$ fermion flavors is always
suppressed. Our gamma matrices are Hermitian and for the trace over the spinor
index we set $\tr 1=4$ for any~$D$. We also set
\begin{equation}
   D=4-2\epsilon.
\label{eq:(2.4)}
\end{equation}

Assuming the dimensional regularization, one can derive a Ward--Takahashi
relation associated with the translational invariance straightforwardly. We
consider the following infinitesimal variation of integration variables in the
functional integral:
\begin{equation}
   \delta A_\mu(x)=\xi_\nu(x)F_{\nu\mu}(x),\qquad
   \delta\psi(x)=\xi_\mu(x)D_\mu\psi(x).
\label{eq:(2.5)}
\end{equation}
Then, since the action changes as\footnote{Here, to make the perturbation
theory well defined, we implicitly assume the existence of the gauge-fixing
term and the Faddeev--Popov ghost fields. However, since they do not explicitly
appear in correlation functions of gauge-invariant operators, we neglect these
elements in what follows.}
\begin{equation}
   \delta S=-\int\mathrm{d}^Dx\,\xi_\nu(x)\partial_\mu T_{\mu\nu}(x),
\label{eq:(2.6)}
\end{equation}
where the energy--momentum tensor~$T_{\mu\nu}(x)$ is defined by
\begin{align}
   T_{\mu\nu}(x)
   &\equiv\frac{1}{g_0^2}\left[
   F_{\mu\rho}^a(x)F_{\nu\rho}^a(x)
   -\frac{1}{4}\delta_{\mu\nu}F_{\rho\sigma}^a(x)F_{\rho\sigma}^a(x)
   \right]
\notag\\
   &\qquad{}
   +\frac{1}{4}
   \Bar{\psi}(x)\left(\gamma_\mu\overleftrightarrow{D}_\nu
   +\gamma_\nu\overleftrightarrow{D}_\mu\right)\psi(x)
   -\delta_{\mu\nu}\Bar{\psi}(x)
   \left(\frac{1}{2}\overleftrightarrow{\Slash{D}}
   +m_0\right)\psi(x),
\label{eq:(2.7)}
\end{align}
with
\begin{equation}
   \overleftrightarrow{D}_\mu\equiv D_\mu-\overleftarrow{D}_\mu,\qquad
   \overleftarrow{D}_\mu\equiv\overleftarrow{\partial}_\mu-A_\mu,
\label{eq:(2.8)}
\end{equation}
we have
% \begin{equation}
%    \int\mathrm{d}^Dx\,\left\langle\partial_\mu T_{\mu\nu}(x)\mathcal{O}\right\rangle
%    =-\left\langle\partial_\nu\mathcal{O}\right\rangle,
% \label{eq:(2.9)}
% \end{equation}
\begin{equation}
   \left\langle\mathcal{O}_{\text{out}}
   \int_{\mathcal{D}}
   \mathrm{d}^Dx\,\partial_\mu T_{\mu\nu}(x)\,\mathcal{O}_{\text{in}}\right\rangle
   =-\left\langle\mathcal{O}_{\text{out}}\,\partial_\nu\mathcal{O}_{\text{in}}
   \right\rangle,
\label{eq:(2.9)}
\end{equation}
where $\mathcal{O}_{\text{in}}$ ($\mathcal{O}_{\text{out}}$) is a collection of
gauge-invariant operators localized inside (outside) the finite integration
region~$\mathcal{D}$. This relation shows that the energy--momentum tensor
generates the infinitesimal translation and, at the same time, the bare
quantity~\eqref{eq:(2.7)} does not receive the multiplicative renormalization.
Thus, we define a renormalized finite energy--momentum tensor by subtracting
its (possibly divergent) vacuum expectation value:
\begin{equation}
   \left\{T_{\mu\nu}\right\}_R(x)
   \equiv T_{\mu\nu}(x)-\left\langle T_{\mu\nu}(x)\right\rangle.
\label{eq:(2.10)}
\end{equation}

A fundamental property of the energy--momentum tensor is the trace
anomaly~\cite{Crewther:1972kn,Chanowitz:1972vd,Adler:1976zt,Nielsen:1977sy,%
Collins:1976yq}. One simple way to derive
this~\cite{Fujikawa:1980rc,Fujikawa:2004cx} is to set $\xi_\mu(x)\propto x_\mu$
in~Eq.~\eqref{eq:(2.5)} and compare the resulting relation with the
renormalization group equation. After some consideration, this yields
\begin{equation}
   \delta_{\mu\nu}
   \left\{T_{\mu\nu}\right\}_R(x)
   =-\frac{\beta}{2g^3}\left\{F_{\rho\sigma}^aF_{\rho\sigma}^a\right\}_R(x)
   -(1+\gamma_m)m\left\{\Bar\psi\psi\right\}_R(x),
\label{eq:(2.11)}
\end{equation}
where we assume that the renormalized operators in the right-hand side are
defined in the $\text{MS}$ scheme~\cite{Collins:1984xc}.\footnote{The
renormalization in the $\text{MS}$ scheme to the one-loop order is summarized
in~Appendix~\ref{sec:A}.} Throughout the present paper, we always assume that
the vacuum expectation value is subtracted in renormalized operators.
In~Eq.~\eqref{eq:(2.11)}, the renormalization group functions $\beta$
and~$\gamma_m$ are defined by
\begin{align}
   \beta&\equiv
   \left(\mu\frac{\partial}{\partial\mu}\right)_0g
   =-\frac{1}{2}g\left(\mu\frac{\partial}{\partial\mu}\right)_0\ln Z,
\label{eq:(2.12)}\\
   \gamma_m&\equiv
   -\left(\mu\frac{\partial}{\partial\mu}\right)_0\ln m
   =-\beta\frac{\partial}{\partial g}\ln Z_m,
\label{eq:(2.13)}
\end{align}
where $g$ and~$m$ are the renormalized gauge coupling and the renormalized
mass, respectively, and the derivative with respect to the renormalization
scale~$\mu$ is taken with all bare quantities kept fixed. The
renormalization constants are defined by
\begin{equation}
   g_0^2=\mu^{2\epsilon}g^2Z,\qquad m_0=mZ_m^{-1}.
\label{eq:(2.14)}
\end{equation}
The first few terms of the perturbative expansion of those renormalization
functions read
\begin{align}
   \beta=-b_0g^3-b_1g^5+O(g^7),\qquad\gamma_m=d_0g^2+d_1g^4+O(g^6),
\label{eq:(2.15)}
\end{align}
where~\cite{Caswell:1974gg,Jones:1974mm}
\begin{align}
   b_0&=\frac{1}{(4\pi)^2}
   \left[\frac{11}{3}C_2(G)-\frac{4}{3}T(R)N_{\mathrm{f}}\right],
\label{eq:(2.16)}\\
   b_1&=\frac{1}{(4\pi)^4}
   \left\{
   \frac{34}{3}C_2(G)^2-\left[4C_2(R)+\frac{20}{3}C_2(G)\right]T(R)N_{\mathrm{f}}
   \right\},
\label{eq:(2.17)}
\end{align}
and~\cite{Tarrach:1980up,Nachtmann:1981zg}
\begin{align}
   d_0&=\frac{1}{(4\pi)^2}6C_2(R),
\label{eq:(2.18)}\\
   d_1&=\frac{1}{(4\pi)^4}
   \left\{
   3C_2(R)^2
   +\left[\frac{97}{3}C_2(G)-\frac{20}{3}T(R)N_{\mathrm{f}}\right]C_2(R)
   \right\}.
\label{eq:(2.19)}
\end{align}
